\newpage
\section{Methodik}\label{methodik}
% Überblick über den gesamten Workflow, idealerweise mit einer
% Abbildung die alle Schritte auf einen Blick zeigt.

% \begin{figure}[H]
% \caption{Methodischer Workflow der Arbeit}\label{fig:workflow}
% \includegraphics[width=0.9\textwidth]{workflow}
% \\
% Quelle: Eigene Darstellung
% \end{figure}

% ============================================================
% Die Gliederung folgt dem Datenfluss von der Befliegung bis zur
% Kennzahl, nicht der zeitlichen Abfolge der Bearbeitung.
% ============================================================

Die folgenden Abschnitte beschreiben den Weg von den bereitgestellten Punktwolken bis zum berechneten Verbuschungsgrad. Die Gliederung folgt dabei dem Datenfluss und nicht der zeitlichen Abfolge der Bearbeitung.

\subsection{Untersuchungsgebiete und Datengrundlage}\label{gebiete_daten}

\subsubsection{Untersuchungsgebiete}\label{untersuchungsgebiet}

Der Authausener Wald liegt in Sachsen nordöstlich von Leipzig und umfasst rund 2.200 Hektar. Ab 1937 beanspruchte das Militär die Fläche, ab 1955 bewirtschaftete sie die DDR als Militärforst und bis 1995 unterhielt die Bundeswehr dort ein Materiallager. Aus dieser Nutzung sind drei kleine Heideflächen innerhalb eines ausgedehnten Waldgebietes hervorgegangen, deren Erhalt und Verbesserung zu den Entwicklungszielen der Fläche zählt \cite{DBU.Authausen}. Beflogen wurden für diese Arbeit vier Gebiete, nämlich die Schneise I, die zusammen beflogenen Schneisen IV und V, die HeideWildnis sowie der Korridor einer 110-kV-Trasse. Bis auf die kompaktere HeideWildnis-Fläche handelt es sich um schmale, von Wald gesäumte Streifen.

Die Cuxhavener Küstenheiden in Niedersachsen umfassen rund 1.400 Hektar und bilden das größte zusammenhängende Heidegebiet auf dem Festland der deutschen Nordseeküste. Auch dort geht die Offenhaltung auf militärische Nutzung zurück, zunächst auf einen Marineschießplatz ab 1915 und später auf einen Truppenübungsplatz der Bundeswehr von 1956 bis 2003. Heute wird die Fläche durch Beweidung mit Schafen, Pferden und Rindern offengehalten \cite{DBU.Cuxhaven}.

\subsubsection{Befliegung und photogrammetrische Auswertung}\label{datengrundlage}

Befliegung und photogrammetrische Auswertung lagen nicht in der Hand dieser Arbeit, sondern erfolgten im Rahmen des Projekts durch die Deutsche Bundesstiftung Umwelt, die die fertigen Orthomosaike und Punktwolken bereitgestellt hat. Die folgenden Angaben geben deshalb wieder, was seitens der DBU dokumentiert vorliegt.

Zum Einsatz kam eine DJI Mavic 3M. Sie trägt eine RGB-Kamera mit 20 Megapixeln sowie vier Multispektralkameras mit je fünf Megapixeln, darunter eine für das Nahinfrarot bei 860 Nanometern. Aus diesem Band stammt der Nahinfrarotkanal, dessen Nutzen
als zusätzliches Eingabemerkmal in Abschnitt~\ref{ablation} geprüft wird. Georeferenziert wurde über das RTK-Modul der Drohne, ohne dass Bodenpasspunkte zum Einsatz kamen. Für den fixierten RTK-Betrieb gibt der Hersteller eine Lagegenauigkeit von einem Zentimeter zuzüglich eines Millionstels der Basislinienlänge an \cite{DJI.Mavic3M}.

Die Authausener Flächen wurden im Mai 2025 beflogen, die Cuxhavener Küstenheide ebenfalls im Jahr 2025. Der Aufnahmezeitpunkt ist insofern von Bedeutung, als sich das Erscheinungsbild der Heide im Jahresverlauf deutlich verändert und blühende Besenheide spektral anders auftritt als nicht blühende.

Die Rekonstruktion der Punktwolken erfolgte mit Agisoft Metashape. Aus ihr gingen fünf getrennte Punktwolken hervor, vier für den Authausener Wald mit zusammen 100,6 Millionen Punkten und eine für die Cuxhavener Küstenheide mit 415,3 Millionen Punkten. Annotierte Teilflächen liegen allein in der Schneise I, in der 110-kV-Trasse und in Cuxhaven, sodass sich Training und die Bewertung anhand gelabelter Testkacheln auf diese drei Wolken stützen, während die übrigen beiden erst bei der Anwendung des Modells auf die gesamte Fläche hinzukommen (Abschnitt~\ref{inferenz_preprocessing}). Die Punktdichte liegt in den Authausener Gebieten zwischen 116 und 156 Punkten je Quadratmeter und in Cuxhaven bei rund 98. Sie bestimmt mit, wie viele Punkte dem Netz bei einem gegebenen Nachbarschaftsradius zur Verfügung stehen. Flughöhe, Bodenauflösung, Überlappungsgrad sowie die verwendete Programmversion und ihre Einstellungen sind dagegen nicht dokumentiert, was für diese Arbeit unerheblich ist, weil die Punktwolken ohnehin als gegebene Eingabe verarbeitet werden.

\subsubsection{Manuelle Annotation und Klassenschema}\label{annotation}

Die Grundlage für Training und Bewertung bilden manuell annotierte Referenzflächen, für die einzelne Ausschnitte der Orthomosaike bearbeitet wurden, im Folgenden Teilflächen genannt. Statt die Flächen von Hand zu zeichnen, zerlegt ein Segmentierungsverfahren jedes Orthomosaik zunächst in kleine, in sich einheitliche Bildregionen, die anschließend manuell einer Klasse zugeordnet werden. Verwendet wurde dafür in R das Verfahren Simple Non-Iterative Clustering (SNIC), das benachbarte Pixel nach spektraler Ähnlichkeit und räumlicher Nähe zu kompakten, zusammenhängenden Regionen zusammenfasst \cite{Achanta.2017}. Die daraus entstehenden Segmente liegen als Polygone vor.

Vergeben wurden die acht Klassen \texttt{bush}, \texttt{deadwood},
\texttt{graminoid}, \texttt{heath}, \texttt{other}, \texttt{sand},
\texttt{soil} und \texttt{tree}. Für die Fragestellung tragen vor allem \texttt{bush} und \texttt{tree} bei, die zusammen die Gehölze bilden, sowie \texttt{heath} und \texttt{graminoid} als die davon abzugrenzende niedrige Vegetation. Eine neunte Klasse für Wasser ist in der Verarbeitung vorgesehen, in den Annotationen jedoch nicht belegt.

Neben der Klasse wurden je Segment drei Angaben zur Aufnahmequalität erfasst. Zwei davon sind binär und halten fest, ob ein Segment verschattet oder überbelichtet ist. Die dritte gibt die Sicherheit der Zuordnung auf einer Skala von null bis zehn an. Diese Angaben steuern das Gewicht eines Punktes im Verlust (Abschnitt~\ref{verlustfunktion}).

Annotiert wurden 16 Teilflächen mit 76.193 Segmenten auf 8,92 Hektar, elf davon im Authausener Wald mit 3,39 Hektar und fünf in der Cuxhavener Küstenheide mit 5,53 Hektar.

\subsection{Datenvorverarbeitung}\label{datenvorverarbeitung}

Die Vorverarbeitung überführt drei Ausgangsgrößen in einen Trainingsdatensatz für PointNet++, nämlich die SfM-Punktwolken der drei Gebiete, die manuell annotierten Polygone und die daran hängenden Qualitätsattribute. Sie ist vollständig in R umgesetzt und stützt sich auf das Paket \texttt{lidR} für die Verarbeitung der Punktwolken \cite{Roussel.2020} sowie auf \texttt{sf} für die Vektoroperationen \cite{Pebesma.2018} und \texttt{terra} für die Rasteroperationen \cite{Hijmans.2025}. Am Ende steht ein Verzeichnis aus Textdateien, in dem jede Datei eine Kachel der Fläche beschreibt und jede Zeile einen Punkt mit seinen Merkmalen und seinem Label enthält.

\subsubsection{Aufbereitung und Prüfung der Ausgangsdaten}\label{georeferenz}

Voraussetzung für den Label-Transfer ist, dass die annotierten Polygone und die Punktwolke im selben Koordinatensystem liegen und räumlich deckungsgleich sind. Beides wurde vorab geprüft, und zwar zuerst rechnerisch über einen Vergleich der EPSG-Codes und der Ausdehnungen und anschließend visuell durch eine gemeinsame Darstellung von Punkten und Polygonrändern. Alle verwendeten Daten liegen in der UTM-Zone 32N (EPSG 25832) vor.

Geladen wird je Teilfläche nur der Ausschnitt, der innerhalb des umschließenden Rechtecks (bounding box) der zugehörigen Polygone liegt. Da diese Rechtecke die tatsächlichen Annotationsflächen nur grob umschreiben, enthalten die geladenen Ausschnitte auch Punkte außerhalb der Polygone, die im weiteren Verlauf als ungelabelt aussortiert werden. Von den mitgeführten Punktattributen werden ausschließlich die Koordinaten, die vier Farbkanäle sowie die Klassen- und Qualitätsangaben übernommen, da die Punktwolken der drei Gebiete unterschiedliche Zusatzattribute enthalten und sich sonst nicht zu einem gemeinsamen Datensatz zusammenführen ließen.

\subsubsection{Bodenerkennung und Höhennormalisierung}\label{hoehennormalisierung}

Die Z-Koordinate der Rohpunkte gibt die absolute Höhe an und ist für eine vegetationskundliche Auswertung wenig aussagekräftig, weil sie das Relief des Geländes enthält. Für die Unterscheidung von Heide, Gras und Gehölzen ist dagegen die Höhe über dem Boden entscheidend, das heißt die Wuchshöhe der Pflanze an ihrem Standort. Die Punktwolke wird deshalb höhennormalisiert, wodurch jeder Punkt statt seiner absoluten Höhe seinen vertikalen Abstand zur Geländeoberfläche erhält.

Zunächst werden die Bodenpunkte automatisch mit einem Cloth Simulation Filter (CSF) erkannt \cite{Zhang.2016}. Anschaulich lässt sich das Verfahren so verstehen, dass die Punktwolke auf den Kopf gestellt wird und ein virtuelles Tuch auf die dabei entstandene Oberfläche fällt, dessen Form sich nach dem Absinken der Geländeoberfläche annähert. Alle Punkte, deren Abstand zu dieser Fläche unter einem Schwellenwert liegt, werden als Boden klassifiziert. Gegenüber älteren hangneigungs- und morphologiebasierten Filtern kommt CSF mit wenigen, einfach zu setzenden Parametern aus \cite{Zhang.2016}. Gesetzt wurden der horizontale Abstand zweier benachbarter Tuchknoten (grid resolution) auf 0,5 Meter, der Abstandsschwellenwert für die Trennung von Boden- und Objektpunkten auf 0,3 Meter sowie die Steifigkeit des Tuchs (rigidness) auf zwei, den mittleren der drei zulässigen Werte.

Aus den so bestimmten Bodenpunkten wird anschließend für jeden Punkt der Wolke die lokale Geländehöhe geschätzt und von seiner Z-Koordinate abgezogen. Die Schätzung erfolgt über eine inverse Distanzgewichtung (inverse distance weighting) der zehn nächstgelegenen Bodenpunkte, deren Beitrag zur Schätzung mit dem Quadrat ihres Abstands zum geschätzten Punkt abnimmt. Bodenpunkte liegen danach nahe null und die Höhe eines Vegetationspunktes entspricht seiner Wuchshöhe.

Diese Normalisierung erfolgt getrennt für jede der 16 annotierten Teilflächen, in deren Zuschnitt die Daten bereits geladen wurden. An den Rändern der Gebiete, wo die photogrammetrische Rekonstruktion ungenau wird und das Geländemodell aus weiter entfernt liegenden Bodenpunkten extrapoliert werden muss, entstehen dadurch stark negative Höhen. Aus diesem Grund werden nach der Normalisierung Punkte unterhalb von minus einem Meter verworfen.

\subsubsection{Übertragung der Labels auf die Punktwolke}\label{label_transfer}

Auf der so normalisierten Punktwolke lassen sich die Annotationen übertragen, die als zweidimensionale Polygone im Luftbild vorliegen, während das Modell ein Label je Punkt der dreidimensionalen Wolke benötigt. Ein räumlicher Verschnitt prüft dafür für jeden Punkt, in welchem Polygon seine X- und Y-Koordinate liegt, und weist ihm dessen Klasse zu. Auf demselben Weg werden neben der Vegetationsklasse auch die drei Qualitätsattribute \texttt{shadow}, \texttt{overexposed} und \texttt{confidence} übertragen. Punkte ohne zugeordnetes Polygon werden anschließend entfernt.

Dieses Vorgehen beantwortet die erste Teilfrage nach einem automatisierten Label-Transfer, bringt jedoch eine systematische Schwäche mit sich. Die
Polygone wurden auf der von oben sichtbaren Oberfläche gezeichnet, ihre Klasse
wird beim Verschnitt aber auf die gesamte darunterliegende Punktsäule
übertragen. Jedes Polygon wirkt damit wie ein Ausstecher, der seine Klasse
durch die volle Höhe der Punktwolke stempelt.

Ragt eine Baumkrone über eine Heidefläche, ist im Luftbild nur die Krone zu sehen, wodurch die darunter wachsende Heide fälschlich das Label \texttt{tree} erhält. An den Rändern derselben Krone kehrt sich der Fehler um, weil einzelne Zweige über die benachbarte Fläche hinausreichen, im Bild aber nicht eigens abgegrenzt werden, da die Polygongrenze vereinfacht entlang der sichtbaren Kronenkante verläuft. Punkte auf diesen Zweigen fallen dadurch in ein Heide- oder Graspolygon und tragen dessen Label. Der Ausstecher wirkt damit in beide Richtungen.

Ein erheblicher Teil der als \texttt{tree} gelabelten Punktmasse liegt, wie die Höhenverteilung je Klasse zeigt, nahe der Bodenoberfläche und bildet damit Unterwuchs statt Baum ab. In der Gegenrichtung reichen einzelne als \texttt{heath} gelabelte Punkte bis 22 Meter, als \texttt{graminoid} gelabelte bis 23 Meter und als \texttt{sand} gelabelte bis 14 Meter über den Boden, was in der Natur für diese Klassen ausgeschlossen ist. Betroffen sind davon allerdings nur wenige Punkte, denn das 95. Perzentil der Heidepunkte liegt bei 1,26 Metern. Anschaulich wird das durch eine Abbildung, die einen senkrechten Schnitt durch die Punktwolke darstellt (Abbildung~\ref{fig:ausstecher}), in dem beide Fehlerarten nebeneinander auftreten.

\begin{figure}[H]
    \centering
    \includegraphics[width=\textwidth]{ausstecher_problem}
    \caption[Fehllabel durch den Label-Transfer]{Senkrechter Schnitt durch die
    Punktwolke der Schneise I nach dem Label-Transfer. Eigene Darstellung.}
    \label{fig:ausstecher}
\end{figure}

\subsubsection{Höhenbasierte Labelkorrektur}\label{labelkorrektur}

Um diese Fehllabel zu bereinigen, wird auf der normalisierten Höhe ein Regelwerk angewendet, das Punkte in zwei Richtungen umklassifiziert. Die Aufwärtsregeln stufen Punkte, die für ihre Klasse zu hoch liegen, zur Klasse \texttt{tree} hoch. Die Abwärtsregeln behandeln den umgekehrten Fall und weisen tiefliegende Punkte unter Baumpolygonen dem Unterwuchs zu, wobei Punkte unterhalb von 0,5 Metern die Klasse \texttt{heath} und Punkte zwischen 0,5 und 2,0 Metern die Klasse \texttt{bush} erhalten (Abbildung~\ref{fig:labelkorrektur}). Die Rückstufung ist dadurch gerechtfertigt, dass die Punktwolke photogrammetrisch aus Aufnahmen von oben entstanden ist und Baumstämme unter geschlossenen Kronen kaum rekonstruiert werden, sodass tiefliegende Punkte unter einem Baumpolygon mit hoher Wahrscheinlichkeit Vegetation am Boden abbilden.

\begin{figure}[H]
    \centering
    \includegraphics[width=\textwidth]{labelkorrektur_zoom}
    \caption[Wirkung der höhenbasierten Labelkorrektur]{Senkrechter Schnitt
    durch die Punktwolke der 110-kV-Trasse vor (a) und nach (b) der
    Labelkorrektur. Die Linien bei 0,5 und 2,0 Metern markieren die Grenzen der Abwärtsregeln.
    Eigene Darstellung.}
    \label{fig:labelkorrektur}
\end{figure}

Um die Schwellenwerte der Aufwärtsregeln aus den Daten selbst abzuleiten, wurde für jede Klasse über alle 16 annotierten Teilflächen das 95. Perzentil der normierten Höhe berechnet und auf einen gerundeten Wert festgelegt, woraus sich ein Höhenschwellenwert von 1,15 Meter für \texttt{soil}, 1,4 Meter für \texttt{sand}, 1,2 Meter für \texttt{graminoid} und 1,3 Meter für \texttt{heath} ergibt.

Auf die Klasse \texttt{bush} ist dieses Vorgehen nicht anwendbar, da ihr oberes Perzentil durch überstehende Bäume so stark verfälscht wird, dass es mit 13,5 Metern weit oberhalb jeder realistischen Strauchhöhe liegt. Ihre Schwelle wurde deshalb als eigene Setzung auf 3,5 Meter festgelegt. Dieser Wert orientiert sich an der Wuchshöhe der auf Heideflächen typischen Pioniergehölze wie Schlehe, Weißdorn und Rosenarten \cite{Schoknecht.1998} und liegt bewusst unterhalb der in der Vegetationskunde üblichen Obergrenze der Strauchschicht von etwa fünf Metern, da Gehölze dieser Höhe im Untersuchungsgebiet bereits durchweg als Birke oder Kiefer und damit als Baum auftreten.

Die Klasse \texttt{deadwood} bleibt bewusst ohne Höhenregel, weil Totholz sowohl liegend als auch stehend vorkommt und ihre Höhenverteilung deshalb zweigipflig ist, sodass eine einzelne Schwelle eine der beiden Erscheinungsformen systematisch falsch behandeln würde.

Die Zuweisung des Unterwuchses über feste Höhenbänder bleibt allerdings eine Näherung. Tiefliegende Punkte unter Baumkronen können auch offener Boden oder Sand sein, werden durch die Abwärtsregeln aber je nach Höhe entweder der Heide oder den Büschen zugeordnet. Beide Klassen können dadurch überschätzt werden, was im Fall von \texttt{bush} sogar unmittelbar die Zielgröße dieser Arbeit betrifft.

Hinzu kommt, dass die Regeln punktweise und nicht objektbezogen arbeiten. Ein Strauch, dessen Basis am Boden ansetzt, wird an der Grenze von 0,5 Metern waagerecht durchgeschnitten und unterhalb davon der Heide, oberhalb den Büschen zugeordnet, obwohl es sich um dieselbe Pflanze handelt. In Abbildung~\ref{fig:labelkorrektur} ist das an mehreren Stellen zu erkennen.

Aus der Reihenfolge der Regeln entsteht eine weitere Nebenwirkung, denn weil die Aufwärtsregeln zuerst greifen, wird etwa ein Graspunkt oberhalb von 1,2 Metern zunächst zu \texttt{tree} hochgestuft und unmittelbar danach von der Abwärtsregel zu \texttt{bush} zurückgestuft, obwohl kein Regelsatz diesen Übergang direkt vorsieht. Betroffen sind davon vor allem die Klassen \texttt{graminoid} und \texttt{heath}. Wie stark beide Effekte die Buschklasse verschieben, wird bei der Einordnung der Ergebnisse aufgegriffen (Kapitel~\ref{diskussion}). Die Schwellenwerte selbst beruhen zudem auf den 16 annotierten Teilflächen und wären bei einer erneuten Erweiterung des Datensatzes neu zu bestimmen.

% ZAHLEN FÜR DIE DISKUSSION (vollständige Punktwolke, 11.721.599 Punkte):
% - Umklassifiziert insgesamt: 1.157.068 Punkte (9,9 %)
% - tree->heath 397.888 | tree->bush 378.590 | graminoid->bush 172.379
%   heath->bush 79.678 | bush->tree 47.663 | heath->tree 35.803
%   graminoid->tree 34.517 | soil/sand->bush,tree 10.550
% - Zerschnittene Strukturen: von 91.026 Rasterzellen (0,25 m) mit
%   umklassifiziertem Unterwuchs enthalten 17.131 danach heath UND bush (18,8 %)
% - Zweierkette (erst hoch zu tree, dann runter zu bush): 258.531 Punkte,
%   also 40,6 % aller Punkte, die neu die Klasse bush erhalten

\subsubsection{Lokale Höhenmerkmale}\label{hoehenmerkmale}

Die bisher vorliegenden Merkmale beschreiben jeden Punkt über seine Lage, seine Farbe und seine Höhe über dem Boden. Für die entscheidende Trennung reicht das nicht aus, denn die Höhe allein trennt einen 40 Zentimeter hohen Busch nicht von einem gleich hohen Heidebestand. Im Luftbild erkennt eine annotierende Person den Unterschied an der Textur, also an derselben unruhigen Oberfläche, die sich in der Punktwolke als kleinräumige Höhenvariabilität niederschlägt. Um diese Information für das Netz verfügbar zu machen, werden je Punkt zwei zusätzliche Merkmale berechnet. Das Merkmal \texttt{z\_range} gibt die Höhenspanne innerhalb einer Rasterzelle von 0,5 Metern an und erreicht große Werte bei locker aufgebauten Sträuchern, deren Zweige in verschiedenen Höhen liegen, kleine Werte im geschlossenen Heideteppich und nahezu den Wert null auf offenem Boden. Das Merkmal \texttt{z\_rel} beschreibt die Höhe eines Punktes über der medianen Höhe seiner Umgebung von 2,5 Metern und wird positiv, wenn eine Struktur aus dem umgebenden Bestand herausragt. Berechnet werden beide Merkmale getrennt je Gebiet.

Beide Merkmale werden bei fünf Metern begrenzt, \texttt{z\_rel} zusätzlich nach unten bei minus fünf Metern. Ohne diese Grenze reicht \texttt{z\_range} in Zellen, die eine Baumkrone schneiden, bis über 20 Meter und drückt den für die Verbuschung entscheidenden Bereich zwischen Heide und Busch auf einen sehr schmalen Ausschnitt zusammen. Verloren geht durch die Kappung wenig, da solche Spannen nur unter Baumkronen auftreten und sich Bäume bereits über die normalisierte Höhe zuverlässig von der niedrigen Vegetation abheben.

\subsubsection{Kachelung}\label{tiling}

Nachdem alle Merkmale berechnet sind, muss die Punktwolke in die Form
gebracht werden, die das Netz erwartet, denn PointNet++ verarbeitet keine
beliebig großen Szenen, sondern Punktmengen fester Größe. Ein gleitendes
Fenster (sliding window) von zehn Metern Kantenlänge wandert dafür über die
Fläche und schneidet daraus quadratische Kacheln. Zwischen zwei Fenstern
liegen jeweils fünf Meter Versatz, sodass benachbarte Kacheln sich zur Hälfte
überlappen. Damit verdoppelt sich nicht nur die Zahl der Trainingsbeispiele, sondern das Netz sieht dieselbe Struktur auch an wechselnder Stelle innerhalb der Kachel, was einer Datenerweiterung (data augmentation) entspricht.

Das gleitende Fenster wird je zusammenhängendem Bereich verschoben und nicht über alle Gebiete hinweg, damit die großen unannotierten Zwischenräume ausgelassen werden. Die acht aneinandergrenzenden Teilflächen der Schneise I bilden ein durchgehendes Raster, wohingegen alle übrigen Teilflächen isoliert liegen und einzeln verarbeitet werden.

Jede Kachel muss dieselbe Anzahl an Punkten enthalten, weil das Netz eine feste Eingabegröße erwartet. Aus jedem Fenster werden deshalb 8.192 Punkte ohne Zurücklegen gezogen, damit kein Punkt doppelt in eine Kachel gelangt und dort eine Dichte vortäuscht, die in den Daten nicht vorhanden ist. Enthält ein Fenster weniger Punkte, wird es verworfen, was vor allem Randfenster betrifft, die nur zum Teil innerhalb einer annotierten Fläche liegen. Aus 4.170 Fenstern entstehen auf diesem Weg 3.017 Kacheln. Gezogen wird mit einem festen Startwert des Zufallsgenerators, damit ein erneuter Export dieselben Kacheln erzeugt und die Läufe der Ablationsstudie vergleichbar bleiben. Wie viele Punkte eine Kachel enthält, gehört zu den Faktoren, die dort variiert werden (Abschnitt~\ref{ablation}).

Die Punkte jeder Kachel werden auf die Fenstermitte zentriert, damit das Netz keine absoluten Koordinaten lernt, sondern nur die Anordnung der Punkte zueinander. Die Farbkanäle einschließlich des Nahinfrarotkanals erhalten Werte zwischen null und eins. Die Klassenbezeichner werden von der Nummerierung eins bis acht auf null bis sieben umgeschrieben, so wie es die Implementierung erwartet. Abgelegt wird jede Kachel als Textdatei, in der eine Zeile einem Punkt entspricht. Dort stehen in fester Reihenfolge die zentrierten X- und Y-Koordinaten, die normalisierte Höhe, die vier Farbkanäle, die beiden lokalen Höhenmerkmale, das Label sowie die drei Qualitätsattribute.

\subsubsection{Räumlich getrennte Datenaufteilung}\label{split}

Bevor trainiert werden kann, muss der Kacheldatensatz in einen Trainings-, einen Validierungs- und einen Testteil zerlegt werden, damit die abschließende Bewertung auf Kacheln beruht, die das Netz nie gesehen hat.

Die Kacheln zufällig auf die drei Teile zu verteilen, wäre in diesem Fall irreführend, denn wegen der 50-prozentigen Überlappung enthalten zwei benachbarte Kacheln zur Hälfte dieselben Punkte. Landet nun die eine im Training und die andere in der Validierung, wird das Netz an Punkten geprüft, mit denen es trainiert wurde, sodass die gemessene Güte zu optimistisch ausfiele.

Verteilt werden deshalb nicht einzelne Kacheln, sondern räumlich voneinander getrennte Blöcke. Alle Kacheln eines Blocks gehen geschlossen in denselben Teil, sodass keine Überlappung eine Grenze überschreiten kann. Die isoliert liegenden Teilflächen bilden jeweils einen eigenen Block, wohingegen die Schneise I als großes, zusammenhängendes Gebiet in drei Blöcke aufgeteilt wird, zwischen denen jeweils ein zehn Meter breiter Streifen frei bleibt, sodass kein Fenster über eine Blockgrenze reichen kann. Ohne diese Aufteilung fiele der gesamte Landschaftstyp der Schneise I in einen einzigen Teil, sodass das Netz auf ihm entweder nicht trainiert oder nicht geprüft würde.

Verteilt werden die Blöcke mit dem Ziel einer Aufteilung von 70 zu 15 zu 15 Prozent, die sich nur annähern lässt, weil immer ganze Blöcke vergeben werden. Von den 3.017 Kacheln entfallen am Ende 2.094 auf das Training, 455 auf die Validierung und 468 auf den Test, was 69, 15 und 16 Prozent entspricht. Beide Untersuchungsgebiete, der Authausener Wald und die Cuxhavener Küstenheide, sind in allen drei Teilen vertreten, sodass keiner der beiden Landschaftstypen ausschließlich im Training oder ausschließlich im Test vorkommt.

\subsection{Modell und Trainingssetup}\label{modell_training}

\subsubsection{Architektur}\label{modell_pointnetpp}
Als Netz kommt PointNet++ in der Variante für semantische Segmentierung zum Einsatz, deren Aufbau in Abschnitt~\ref{pointnet++} beschrieben ist \cite{Qi.2017b}. Gewählt wurde es, weil es Punktwolken unmittelbar verarbeitet und die Höhenstruktur des Bestandes erhält, während zweidimensionale Verfahren die dritte Dimension auf ein Rasterbild projizieren müssen. Genau an dieser Struktur hängt die Unterscheidung zwischen Busch und Heide.

Die Umsetzung folgt der frei verfügbaren Referenzimplementierung von Yan \cite{Yan.2019}, die unter der MIT-Lizenz steht. Übernommen wurde sie unverändert als eigener Commit, bevor die erste eigene Zeile entstand, sodass sich jede spätere Änderung dieser Arbeit zuordnen lässt.

Das Netz betrachtet eine Kachel in vier Set-Abstraction-Stufen mit jeweils größerem Blickfeld. In der ersten Stufe wählt es 1.024 Punkte als Zentren aus und beschreibt jedes davon anhand von bis zu 32 Nachbarn im Umkreis von 0,1 Metern, wodurch feine Textur erfasst wird. Über die folgenden Stufen sinkt die Zahl der Zentren auf 256, 64 und schließlich auf 16, während der Umkreis auf 0,2, 0,4 und 0,8 Meter wächst. Die Szene wird dadurch von Stufe zu Stufe gröber betrachtet, dafür aber im größeren Zusammenhang. Vier Feature-Propagation-Stufen verteilen die so gewonnenen Beschreibungen anschließend wieder auf alle Punkte, sodass am Ende jeder einzelne Punkt einen Wert je Klasse erhält. Eine Dropout-Schicht setzt während des Trainings in jedem Schritt zufällig die Hälfte der Merkmalswerte auf null, sodass sich das Netz nicht auf einzelne Merkmale verlassen kann und weniger zur Überanpassung (overfitting) neigt. Bei der Anwendung des fertigen Modells bleibt sie abgeschaltet. Eingabekanäle, Nachbarschaftsradien und die Zahl der Radien je Stufe sind umschaltbar und gehören zu den Faktoren, die in Abschnitt~\ref{ablation} variiert werden.

\subsubsection{Verlustfunktion}\label{verlustfunktion}
Was ein Netz lernt, steuert die Verlustfunktion (loss function). Sie misst, wie weit die Vorhersage vom Label abweicht, und genau dieser Wert wird beim Training minimiert. Zwei Eigenschaften der Trainingsdaten erfordern an dieser Stelle einen Eingriff.

Die erste ist das Klassenungleichgewicht (class imbalance), denn Bäume und Heide stellen zusammen den weit überwiegenden Teil der Punkte, während \texttt{deadwood} und \texttt{soil} nur vereinzelt vorkommen. Ein Netz kann den Fehler am einfachsten senken, indem es die seltenen Klassen gar nicht erst vorhersagt. Um dem entgegenzuwirken, geht jede Klasse mit einem eigenen Klassengewicht in den Verlust ein. Es berechnet sich aus dem Kehrwert der Wurzel ihrer Häufigkeit und fällt damit umso größer aus, je seltener die Klasse ist. Die Wurzel schwächt den Ausgleich bewusst ab, denn ohne sie bekämen die seltensten Klassen ein so hohes Gewicht, dass das Training instabil würde.

Die zweite Eigenschaft ist die unterschiedliche Verlässlichkeit der Labels, denn nicht jedes annotierte Segment war im Luftbild gleich gut zu erkennen. Aus den drei Qualitätsattributen wird deshalb je Punkt ein zusätzliches Gewicht abgeleitet, das unsichere und schlecht sichtbare Punkte abschwächt, ohne sie ganz auszuschließen. In der Praxis wirkt es kaum, denn von den 76.193 Segmenten sind nur 220 als verschattet oder überbelichtet markiert, während 75.916 den höchsten Sicherheitswert tragen.

Beide Gewichtungen greifen nacheinander, denn der Fehler eines Punktes wird zuerst mit dem Gewicht seiner Klasse multipliziert und anschließend über alle Punkte gemittelt, wobei jeder Punkt entsprechend seinem Qualitätsgewicht in dieses Mittel eingeht. Die Qualitätsattribute sind damit kein Eingabemerkmal des Netzes, sondern steuern nur, wie stark ein Punkt beim Lernen zählt.

\subsubsection{Training}\label{training}

Die Trainingsparameter bleiben über alle Läufe unverändert, damit Unterschiede zwischen ihnen allein auf den jeweils geänderten Faktor zurückgehen. Optimiert wird mit dem Adam-Verfahren \cite{Kingma.2015}, das auch die Referenzimplementierung verwendet. Die Lernrate beginnt bei 0,001 und fällt alle zehn Epochen auf 70 Prozent ihres Wertes, sodass das Netz zunächst grob und gegen Ende fein nachjustiert. Hinzu kommt ein Gewichtsverfall (weight decay) von 0,0001. Das Training läuft über 64 Epochen in Stapeln (batches) von 32 Kacheln, wobei es ab run08 nur noch 16 sind, weil die verdoppelte Punktzahl je Kachel den Speicherbedarf entsprechend erhöht. Gerechnet wurde in einer Conda-Umgebung mit Python 3.9 auf einer CUDA-fähigen Grafikkarte im JupyterHub der Universität.

Neben der Überlappung der Kacheln (Abschnitt~\ref{tiling}) kommt während des Trainings eine zweite Form der Datenerweiterung hinzu. Jede Kachel wird vor der Verarbeitung um die Z-Achse um einen zufälligen Winkel gedreht. Auf ihre Koordinaten wird zusätzlich ein Rauschen mit einer Standardabweichung von einem Zentimeter, begrenzt auf fünf Zentimeter, hinzugerechnet. Die Drehung ist fachlich begründet, denn in welche Himmelsrichtung eine Kachel zeigt, sagt nichts über die Vegetation darin aus. Das Rauschen macht das Netz unempfindlicher gegen die Ungenauigkeiten der photogrammetrischen Rekonstruktion.

Gespeichert wird der Zwischenstand mit der höchsten mittleren IoU auf dem Validierungsteil. Als Kriterium dient bewusst nicht der Verlust, sondern die mittlere IoU, weil sie alle Klassen gleich gewichtet und damit derselben Größe folgt, an der die Ergebnisse später gemessen werden.

\subsection{Ablationsstudie}\label{ablation}

Der Studie ging eine Sondierungsphase voraus, in der drei Läufe auf einer
kleineren Datenbasis mit sechs annotierten Teilflächen entstanden. Aus ihr
stammen die Kachelparameter, die Wahl der Architektur und die Gewichtung der
Punkte über die Qualitätsattribute. Nach der Erweiterung auf 16 Teilflächen
wurden alle Läufe neu und ohne Warmstart gerechnet, sodass die Vorläufe auf
einer anderen Datengrundlage beruhen und nicht berichtet werden.

Für die eigentliche Studie ändert jeder Lauf gegenüber seinem Vergleichslauf
genau einen Faktor und lässt alles Übrige unverändert. Trainiert wird jedes Mal von Grund auf, statt auf einem vorhandenen
Modell aufzusetzen, denn nur so lässt sich eine Veränderung dem geänderten
Faktor zuordnen. Möglich wird dieses Vorgehen dadurch, dass sich Klassenschema
und Nachbarschaftsradien über Umgebungsvariablen umstellen lassen und die Zahl
der Eingabemerkmale aus der Spaltenzahl der Kacheldateien abgeleitet wird.
Trainings- und Auswertungsskript bleiben für alle Varianten unverändert,
ergänzt um ein Protokoll, das festhält, welche Einstellungen ein Lauf verwendet
hat. Streng eingehalten wird das bei allen Läufen außer run08, bei dem die
verdoppelte Punktzahl je Kachel zwei weitere Änderungen nach sich zieht, nämlich
eine halbierte Stapelgröße und den Wegfall des Auffüllens unvollständiger
Fenster.

\begin{table}[H]
    \centering
    \caption[Übersicht der Ablationsläufe]{Übersicht der Ablationsläufe mit dem
    jeweils geänderten Faktor und dem Lauf, gegen den verglichen wird.
    Eigene Darstellung.}
    \label{tab:ablation}
    \begin{tabular}{lll}
    \hline
    Lauf & Geänderter Faktor & Vergleich \\
    \hline
    run01 & Ausgangskonfiguration mit acht Klassen und sechs Kanälen & \\
    run02 & Nahinfrarot als siebter Kanal & run01 \\
    run03 & Focal Loss \cite{Lin.2017} statt gewichtetem Verlust & run02 \\
    run04 & Reduktion auf sechs Klassen & run02 \\
    run05 & Zwei lokale Höhenmerkmale als achter und neunter Kanal & run04 \\
    run06 & Reduktion auf zwei Klassen, Gehölz gegen Rest & run05 \\
    run07 & Nachbarschaftsradien um Faktor 2,5 vergrößert & run05 \\
    run08 & 8.192 statt 4.096 Punkte je Kachel & run07 \\
    run09 & Zwei Radien je Stufe (multi-scale grouping) & run08 \\
    \hline
    \end{tabular}
\end{table}

Zwei der neun Läufe erweitern die Merkmale je Punkt, einer verändert die
Verlustfunktion, zwei vereinfachen das Klassenschema und die letzten drei
verändern die Geometrie der Nachbarschaftssuche im Netz, wie Tabelle~\ref{tab:ablation}
zeigt. Diese letzte Gruppe geht auf eine Zwischenauswertung zurück, denn nachdem
mehrere Änderungen an Merkmalen und Verlustfunktion ohne Wirkung geblieben
waren, wurde geprüft, ob die Nachbarschaftsparameter des Netzes überhaupt zur
Punktdichte der Daten passen.
Die Radien der Referenzimplementierung sind auf den Innenraumdatensatz S3DIS
abgestimmt, dessen Blöcke einen Meter Kantenlänge haben und ebenfalls 4.096
Punkte enthalten. Auf Kacheln von zehn Metern übertragen findet die erste
Set-Abstraction-Stufe bei einem Radius von 0,1 Metern im Median nur einen
einzigen Punkt, nämlich ihr eigenes Zentrum, und lernt dort folglich nichts.
Daraus ergaben sich zwei Wege, entweder die Radien zu vergrößern oder die
Punktzahl je Kachel zu erhöhen.

Aus dem Export für run08 ergibt sich eine Einschränkung des Vergleichs, denn
mit ihm änderten sich nicht nur die Trainings-, sondern auch die Testkacheln.
Randfenster mit zu wenigen Punkten blieben bis dahin erhalten, weil sie durch
Kopien vorhandener Punkte aufgefüllt wurden. Nun dagegen werden sie verworfen. Übrig bleiben
dadurch 468 statt 487 Testkacheln. Da Läufe auf unterschiedlichen
Testmengen nicht vergleichbar sind, wurden run05 bis run07 nachträglich gegen
die neue Testmenge ausgewertet. Möglich ist das ohne Neutraining, weil
PointNet++ beliebige Punktzahlen verarbeitet und run05 bis run07 dieselben neun
Eingabekanäle nutzen wie run08 und run09. Die Läufe run05 bis run09 stehen damit
alle auf derselben Testmenge von 468 Kacheln. Für run01 bis run04 gilt das
nicht, denn ihre Eingabeschicht ist auf sechs oder sieben Kanäle festgelegt und
passt damit nicht zu den neun Kanälen des neuen Exports. Diese vier Läufe stehen
deshalb in einer eigenen Tabelle mit der ursprünglichen Testmenge.

\subsection{Evaluation}\label{evaluation}

Bewertet wird jeder Lauf auf dem Testteil, der bis dahin unberührt geblieben
ist. Ein eigenes Skript sammelt die Vorhersagen über alle Testkacheln in einer
gemeinsamen Konfusionsmatrix und leitet daraus die in
Abschnitt~\ref{evaluationsmetriken} definierten Kennzahlen der Segmentierung ab.
Ergänzend gibt das Skript die zehn häufigsten Verwechslungen aus, woran sich
ablesen lässt, zwischen welchen Klassen die Fehler tatsächlich auftreten.

Für die Fragestellung reicht die mittlere IoU allerdings nicht aus, weil sie
über alle Klassen mittelt und deshalb zwischen Läufen mit acht, sechs oder zwei
Klassen nicht vergleichbar ist. Zusätzlich werden die Vorhersagen jedes Modells
nachträglich zu Gehölz, das heißt \texttt{bush} und \texttt{tree}, sowie dem Rest
zusammengefasst und für diese Zweiteilung IoU, Precision, Recall und F1-Score
berechnet. Diese Größe hängt nicht vom trainierten Klassenschema ab und
entspricht unmittelbar dem, was für den Verbuschungsgrad gebraucht wird.
Zugleich ist sie der Maßstab, an dem sich der binäre Lauf messen lassen muss,
denn die Gehölzklasse lässt sich auf zwei Wegen gewinnen, entweder durch ein
eigens auf zwei Klassen trainiertes Modell oder durch nachträgliches
Zusammenfassen der Vorhersagen eines Mehrklassenmodells. Der zweite Weg
verursacht keinen zusätzlichen Aufwand, sodass sich ein eigens binär
trainiertes Modell nur rechtfertigt, wenn es besser abschneidet.

Zu jedem Lauf entsteht ein Protokoll mit Konfusionsmatrix, Kennzahlen je Klasse
und der Zusammenfassung Gehölz gegen Rest, sodass sich die Läufe im Nachhinein
ohne erneute Berechnung gegenüberstellen lassen.

\subsection{Quantifizierung des Verbuschungsgrades}\label{quantifizierung}

Nachdem die vorangegangenen Abschnitte beschrieben haben, wie gut das Netz einzelne Punkte einer Vegetationsklasse zuordnet, geht es nun um die eigentliche Zielgröße der Arbeit. Aus den punktweisen Vorhersagen des in der Ablationsstudie besten Modells, run07, soll ein Verbuschungsgrad je Fläche entstehen, der sich mit einer im Gelände erhobenen Schätzung vergleichen lässt.

\subsubsection{Vorbereitung der Befliegungsgebiete}\label{inferenz_preprocessing}

Während für Training und Bewertung die 16 annotierten Teilflächen ausreichten, muss für den Verbuschungsgrad das vollständige Befliegungsgebiet klassifiziert werden, weil die Flächen der Geländekartierung überwiegend außerhalb der annotierten Ausschnitte liegen. Verarbeitet wurden deshalb die vier Authausener Befliegungen mit zusammen 100,6 Millionen Punkten.

Vorbereitet werden sie mit denselben Schritten und unveränderten Parametern wie die Trainingsdaten (Abschnitte~\ref{hoehennormalisierung} und~\ref{hoehenmerkmale}). Schon eine abweichende Kappungsgrenze oder ein anderer Bodenfilter würde dazu führen, dass derselbe Merkmalswert bei der Anwendung anders zu deuten ist als beim Training.

Die Verarbeitung läuft blockweise in Abschnitten von hundert Metern, denen jeweils ein Puffer von zwanzig Metern beigegeben wird. Eine Bodenklassifikation über das gesamte Gebiet erwies sich als nicht praktikabel, weil Arbeitsspeicherbedarf und Rechenzeit mit der Punktzahl stark ansteigen. Dieser Puffer verhindert Stufen in der Bodeninterpolation an den Blockgrenzen und wird vor dem Schreiben wieder entfernt. Als Ergebnis liegen dann georeferenzierte LAZ-Blöcke vor, die die beiden Höhenmerkmale als zusätzliche Punktattribute mitführen.

\subsubsection{Anwendung des Modells}\label{inferenz}

Während alle bisherigen Läufe der Messung auf gelabelten Kacheln beruhten, wird das Modell nun auf Punkte angewendet, für die keine Annotation vorliegt. Zerlegt werden die Blöcke dabei in Fenster derselben Größe und mit demselben Versatz wie beim Training (Abschnitt~\ref{tiling}), wobei die Fensterursprünge nun nicht mehr am Rand der jeweiligen Fläche ansetzen, sondern auf einem gemeinsamen Raster über alle Blöcke hinweg liegen. Andernfalls wären die Fenster zweier Nachbarblöcke gegeneinander verschoben und ein Punkt an der Grenze fiele in unterschiedlich geschnittene Nachbarschaften. Geladen wird jeder Block zusammen mit einem Saum von zehn Metern aus den angrenzenden Blöcken, vorhergesagt werden jedoch nur seine eigenen Punkte. Ohne diesen Saum sähe das Netz an jeder Blockgrenze nur eine halbe Nachbarschaft.

Aus dem Versatz von fünf Metern bei zehn Metern Kantenlänge (Abschnitt~\ref{tiling}) folgt, dass jeder Punkt in bis zu vier Fenstern liegt und entsprechend mehrfach vorhergesagt wird. Statt diese Vorhersagen auszuzählen und die häufigste zu übernehmen, werden die Klassenwahrscheinlichkeiten über alle beteiligten Fenster gemittelt und erst aus dem Mittel die wahrscheinlichste Klasse bestimmt. Dadurch setzt sich eine sichere Vorhersage gegen eine unsichere durch und Randfehler einzelner Fenster gleichen sich aus. Der gemittelte Höchstwert bleibt als Sicherheitsmaß je Punkt erhalten.

Anders als beim Training, wo unvollständig belegte Fenster verworfen wurden, lässt sich hier nicht auf Bereiche verzichten, für die anschließend ein Verbuschungsgrad berechnet werden soll. Vorhergesagt wird deshalb auch in dünn belegten Fenstern, sofern sie mindestens hundert Punkte enthalten.

\subsubsection{Bezugsfläche und Verbuschungsmaß}\label{verbuschungsmass}

Naheliegend wäre, den Verbuschungsgrad als Anteil der Gehölzpunkte an allen Punkten einer Fläche zu berechnen. Dieses Maß wäre jedoch verzerrt, weil die Punktdichte einer photogrammetrischen Punktwolke ungleich verteilt ist und ein Strauch mit Ästen und Textur je Quadratmeter ein Vielfaches der Punkte liefert, die auf flachem Sandboden entstehen. Unter Gehölzen fehlen zugleich Punkte, weil die Kronen den Boden verdecken, sodass ein Punktanteil die Verbuschung systematisch überschätzen würde.

Verwendet wird stattdessen die projizierte Flächendeckung, also der von oben sichtbare Anteil der Fläche. Dafür wird ein Raster von 0,25 Metern über die Fläche gelegt, in dem jede Zelle Klasse und normalisierte Höhe ihres höchsten vorhergesagten Punktes übernimmt. Das entspricht dem Blick von oben, denn was von oben sichtbar ist, bedeckt den Boden. Ein Heidepunkt unter einem Buschzweig zählt folglich nicht als offene Heide. Da die Referenzdaten auf einer Deckungsgradschätzung im Gelände beruhen, die demselben Prinzip folgt, sind beide Größen unmittelbar vergleichbar.

Als verbuscht gilt eine Zelle dieses Feinrasters, wenn ihre Klasse \texttt{bush} oder \texttt{tree} lautet und ihre Höhe unterhalb einer festgelegten Grenze bleibt. Nötig ist diese zusätzliche Bedingung, weil die Klasse \texttt{tree} zwei fachlich verschiedene Dinge zusammenfasst, nämlich die aufkommenden Pioniergehölze auf der Offenfläche und den angrenzenden Altbestand, dessen Kronen über die Polygongrenze ragen. In der Klasse unterscheiden sich beide nicht, wohl aber in der Höhe. Fachlich deckt sich das mit der Definition, nach der Verbuschung das Aufkommen von Pioniergehölzen auf offener Fläche meint und nicht den benachbarten Wald.

Welchen Wert diese Grenze annimmt, wurde aus dem Vergleich mit der Referenz bestimmt. Berechnet wurden dafür Varianten mit zwei, drei, fünf und acht Metern sowie eine ohne Begrenzung, die nach demselben Verfahren gegen die Kartierung geprüft werden (Abschnitt~\ref{validierung}). Welche Grenze sich durchsetzt, zeigt Abschnitt~\ref{ergebnis_quantifizierung}.

\subsubsection{Referenzdaten und Validierung}\label{validierung}

Als Referenz dient die FFH- und Biotoptypenkartierung des Authausener Waldes aus dem Jahr 2014 \cite{DBU.Biotopkartierung2014}, die 393 Polygone umfasst und im Feld \texttt{VERBUP} den im Gelände geschätzten Verbuschungsgrad in Prozent führt. Ausgewertet werden ausschließlich Offenlandbiotope, Heiden und Magerrasen, Moore und Sümpfe, Staudenfluren sowie Grünland. In Wäldern, Gewässern und Rohbodenbiotopen ist die Größe nicht definiert und steht dort ausnahmslos auf null. Ein Polygon geht nur dann in die Auswertung ein, wenn mindestens 80 Prozent seiner Fläche mit Vorhersagen belegt sind und es mindestens 200 Zellen des Feinrasters umfasst. Somit verbleiben 43 Polygone, von denen 19 auf Heiden und Magerrasen entfallen und 14 von diesen dem Lebensraumtyp 4030, den trockenen europäischen Heiden, zugeordnet sind.

Zwischen Kartierung und Befliegung liegen elf Jahre, was den Vergleich belastet, denn eine 2014 frisch entkusselte Fläche ist bis 2025 wieder aufgewachsen. Um solche Fälle kenntlich zu machen, wird das Bemerkungsfeld der Kartierung nach Hinweisen auf kurz zuvor erfolgte Pflegemaßnahmen durchsucht. Sieben der 43 Polygone tragen einen solchen Vermerk, weshalb alle Kennzahlen zweimal berechnet werden, um eine mögliche Verzerrung des Ergebnisses auszuschließen.

Berichtet werden die in Abschnitt~\ref{evaluationsmetriken} definierten
Kennzahlen des Verbuschungsgrades. Da bei einer Stichprobe dieser Größe
einzelne Flächen die Absolutwerte verzerren können, tritt die Rangkorrelation
nach Spearman hinzu, die allein die Reihenfolge der Flächen bewertet.

Aus demselben Feinraster entsteht schließlich ein Auswerteraster von fünf Metern, dessen Zellen jeweils den Anteil der als Verbuschung vorhergesagten Feinzellen an den kartierten Feinzellen tragen. Zellen, in denen weniger als die Hälfte der Feinzellen belegt ist, bleiben leer, damit dünn abgedeckte Randbereiche keine Scheinwerte liefern. Ausgegeben wird ein GeoTIFF in EPSG~25832, das über die vier Befliegungen 64,5 Hektar mit Werten abdeckt und damit ein Vielfaches der 24,9 Hektar, die die Geländekartierung beisteuert.

Sämtliche Skripte der beschriebenen Verarbeitungskette sind öffentlich verfügbar \cite{Waldeyer.2026}. Das Repository enthält die Vorverarbeitung in R, das Training, die Auswertung und die Berechnung des Verbuschungsgrades, nicht jedoch die Punktwolken und Annotationen.
